% BT1641 — Photonic FSR / Ihara Zeta Bridge
% For inclusion in photonic_holonet.tex
% PASS 5884–5886, 2026-08-17

\section{Spectral Foundation: Ihara Zeta and Free Spectral Range}
\label{sec:ihara_fsr}

The photonic holonet derives its resonance structure directly from the
non-backtracking spectral theory of the W33 graph.

\subsection{Free Spectral Range from Graph Valence}

Let the holonet operate on the W33 twisted-mesh topology with vertex valence $d$.
Define $q = d - 1$. The Ihara zeta function $\zeta_{W_{33}}(u)$ has all its
non-trivial poles on the circle $|u| = q^{-1/2}$ (Ramanujan condition,
Theorem~\ref{thm:equalized_q} of the main paper).
In the physical picture, each pole $u_{k}$ corresponds to a resonant mode;
the pole argument $\phi_k = \arg(1/u_k)$ is the modal phase.

The \emph{free spectral range} is the fundamental period of the pole-argument
sequence:
\[
  \boxed{\mathrm{FSR} = \frac{\log q}{\tau_{\mathrm{rt}}}},
\]
where $\tau_{\mathrm{rt}}$ is the photon round-trip time in the holonet.
For $d = 3$ (W33 is 3-regular in the twisted-mesh model), $q = 2$ and
$\mathrm{FSR} = (\ln 2)/\tau_{\mathrm{rt}} \approx 0.693/\tau_{\mathrm{rt}}$.

\subsection{Decay Rate and Finesse}

The common decay rate (equalized-Q condition) is
\[
  \gamma = \tfrac{1}{2}\log q = \tfrac{1}{2}\ln q.
\]
The resonator finesse, by analogy with the Fabry--P\'erot formula, is
\[
  \mathcal{F} = \frac{\pi\sqrt{q}}{|1-q|}.
\]
For $q = 2$: $\mathcal{F} = \pi\sqrt{2} \approx 4.44$.
This graph-theoretic finesse sets the maximum signal-to-noise ratio achievable
by any routing protocol over the W33 holonet.

\subsection{Mode Equidistribution and Channel Capacity}

By Corollary~\ref{cor:photonic_fsr}, the modal phases are equidistributed
mod FSR. Shannon channel capacity for $N$ equidistributed modes is
\[
  C = N \cdot \log_2\!\left(1 + \frac{\mathcal{F}}{\pi}\right)
  \text{ bits per FSR interval}.
\]
For W33 with $N = 33$ vertices and $\mathcal{F} = \pi\sqrt{2}$:
$C = 33\log_2(1 + \sqrt{2}) \approx 33 \times 1.272 \approx 41.97$ bits/FSR.

\subsection{Hardware Implications}

The FSR and finesse values above are \emph{topology-locked}: they depend only
on the graph valence $d$, not on the physical implementation.
This means:
\begin{itemize}
  \item Any photonic or microwave realization of W33 with valence $d=3$ achieves
    the same finesse $\pi\sqrt{2}$.
  \item The equalized-Q condition ($\gamma$ identical for all modes) implies
    \emph{uniform loss per mode} — a key requirement for fault-tolerant
    operation.
  \item Fabrication tolerances enter only through $\tau_{\mathrm{rt}}$, not
    through $\mathcal{F}$.
\end{itemize}

% Cross-references: BT1640_equalized_q_ihara_theorem.tex,
%   verify_equalized_q_ihara_zeta.py, BREAKTHROUGH_BT679_YANG_MILLS_MASS_GAP.md
